\section{Token 与 ID：文字怎样进入数字世界？}
\label{intro:20a-tokenizer}
\phantomsection\label{concept:token}

上一章已经会用矩阵组织学习，但输入还是数字。现在给模型一句中文，
如何把它变成可以查表的输入？

\subsection{Token、词表、整数 ID 分别做什么？}

\textbf{Token} 是文本处理单位，不一定恰好等于自然语言里的词：
它也可能是单字、子词或标点。
\textbf{Tokenizer（分词器）}按照某套规则把文本拆成 Token，
然后按\textbf{词表}将它们映射为整数 \textbf{ID}。

本书统一使用一份\textbf{人为指定}的五项教学词表：
\[
\begin{array}{c|ccccc}
\text{ID}&0&1&2&3&4\\\hline
\text{Token}&\text{其他}&\text{春天}&\text{。}&\text{花}&\text{开了}
\end{array}
\]
所以“春天 / 花 / 开了 / 。”映射到 $[1,3,4,2]$。
整数编号只是身份，不能因为 ID 4 大于 ID 1 就认为“开了”更重要。

\subsection{子词与 BPE：为什么不一定按汉字切？}

示意短语“今天天气好”如果逐字分割，会成为
\[
[\text{今},\text{天},\text{天},\text{气},\text{好}].
\]
在一个简化的\textbf{子词/BPE 合并演示}里，
若“今+天”和“天+气”经常相邻，可能学到这两次合并：
\[
[\text{今},\text{天},\text{天},\text{气},\text{好}]
\ \longrightarrow\
[\text{今天},\text{天气},\text{好}].
\]
\textbf{这是从汉字出发的教学示意，不是真实 GPT 分词器的逐字节运行记录。}
真正的 BPE 词表及合并规则因模型不同而不同，
中文有时会被切成字、子词甚至字节片段。

切得更细可能增加序列长度，切得较粗可能增加词表体积；
Token 化影响上下文成本，却\textbf{不会自动赋予 ID 语义}。

\begin{warningbox}
本文五项词表只为统一计算和手算服务。
真实模型并不保证将“春天”或“开了”保留为单个 Token。
\end{warningbox}

\par\medskip\noindent{\sffamily\footnotesize\color{blue}\textbf{读完自测}}\quad
\phantomsection\label{check:20a-tokenizer}
按示意规则合并“今+天”“天+气”后，输入的 Token 数从多少变成多少？
\quad{\footnotesize\hyperref[answer:20a-tokenizer]{核对解答。}}\par\medskip

\begin{conclusion}
Token 化将原始文字变成离散 ID，但 ID 只是类别标签。
下一章要解决：如何为每个 ID 学到可用于计算的数值特征？
\end{conclusion}
